%%%%%%%%%%%%%%%%%%%%%%%%%%%%%%%%%%%%%%%%%%%%%%%%%%%%%%%%%%%%%%%%%%%%%%
% Lớp 10 | Đại số tổ hợp | Hoán vị. Chỉnh hợp. Tổ hợp
% Dạng 0D8?2-7: Bài toán xếp chỗ (không tròn, không lặp)
% Số câu: 38  (sinh tự động bởi tach_cau.py; từ nay thầy sửa trực tiếp file này)
%%%%%%%%%%%%%%%%%%%%%%%%%%%%%%%%%%%%%%%%%%%%%%%%%%%%%%%%%%%%%%%%%%%%%%

%%%%%%%%%%%%%%%% TRẮC NGHIỆM 4 PHƯƠNG ÁN %%%%%%%%%%%%%%%%

% nguồn: MaHoaNH_TracNghiem_Lop10
\begin{ex}%[0D8N2-7]%[ma:25470dc9]
	Có bao nhiêu cách sắp xếp $4$ học sinh ngồi vào bàn dài có $8$ chỗ  ngồi?
	 \choice
		 {$\mathrm{C}_{8}^{4}$}
		{\True $\mathrm{A}_{8}^{4}$} 
		{$4!8!$} 
		{$12$} 
	\loigiai{ 
		Mỗi cách xếp $4$ học sinh ngồi vào bàn dài có $8$ chỗ ngồi  chính là một chỉnh hợp chập $4$ của $8$ phần tử.\\
		Vậy có tất cả $\mathrm{A}_{8}^{4}$ cách xếp. 
	} 
\end{ex}

% nguồn: MaHoaNH_TracNghiem_Lop10
\begin{ex}%[0D8N2-7]%[ma:ba0e7417]
	Số cách xếp $5$ bạn An, Bình, Cường, Dũng, Mai vào một dãy ghế  dài có đúng $5$ ghế và mỗi ghế chỉ có đúng một bạn ngồi là
	 \choice
		 {$5$}
		{$36$} 
		{$60$} 
		{\True $120$} 
	 \loigiai{
		Mỗi cách xếp $5$ bạn vào $5$ ghế chính  là một hoán vị của $5$ phần tử.\\
		Do đó có tất cả $5!=120$ (cách  xếp).
	 }
\end{ex}

% nguồn: MaHoaNH_TracNghiem_Lop10
\begin{ex}%[0D8H2-7]%[ma:b23e67ba]
	Có $6 $ học sinh trong đó có bạn An. Hỏi có bao nhiêu cách sắp xếp $6$ học sinh này thành một hàng  dọc sao cho An luôn đứng đầu hàng?
	\choice 
		{\True  $120$}
		{$720$} 
		 {$30$}
		{$24$} 
	 \loigiai{
		Số  cách xếp thỏa yêu cầu bài bằng số hoán vị của $5$ phần tử bằng $5!=120$.
	 }
\end{ex}

% nguồn: MaHoaNH_TracNghiem_Lop10
\begin{ex}%[0D8H2-7]%[ma:366a7c11]
	Có bao nhiêu cách xếp bốn bạn Hà, Mai, Nam, Dương thành một hàng dọc biết bạn Mai đứng đầu  hàng?
	\choice 
		{$7$} 
		 {$24$}
		{\True $6$} 
		 {$12$}
	 \loigiai{
		Vì Mai đứng đầu hàng nên số cách xếp là $1\cdot 3!  = 6$ cách.
	 }
\end{ex}

% nguồn: MaHoaNH_TracNghiem_Lop10
\begin{ex}%[0D8H2-7]%[ma:c9af5801]
	Sắp xếp năm bạn học sinh $A$, $B$, $C$, $D$, $E$ vào một chiếc ghế dài có $5$ chỗ ngồi. Số cách sắp xếp sao cho  bạn $C$ luôn ngồi chính giữa là
	 \choice
		{$120$} 
		{\True $24$} 
		{$60$} 
		 {$16$}
	\loigiai{ 
		Xếp bạn $C$ ngồi giữa có $1$ cách. Mỗi cách xếp  $4$ bạn học sinh $A$, $B$, $D$, $E$ vào $4$ chỗ còn lại là một hoán vị của $4$ phần tử nên có có $4!$ cách. \\
		 Vậy có $24$ cách xếp.
	} 
\end{ex}

% nguồn: MaHoaNH_TracNghiem_Lop10
\begin{ex}%[0D8H2-7]%[ma:95266a50]
	Một nhóm $6$ bạn học sinh mua vé vào rạp chiếu phim. Các bạn mua $6$ vé gồm $3$ vé mang số ghế chẵn, $3$ vé mang số ghế lẻ và không có hai vé nào cùng số. Trong $6$ bạn thì hai bạn muốn ngồi bên ghế chẵn, hai bạn muốn ngồi bên ghế lẻ, hai bạn còn lại không có yêu cầu gì. Hỏi có bao nhiêu cách xếp chỗ  để thỏa mãn các yêu cầu của tất cả các bạn đó?
	 \choice
		{\True $72$} 
		 {$36$}
		 {$18$}
		{$180$} 
	 \loigiai{
		 Số cách chọn $2$ vé cho hai bạn muốn ngồi ghế bên chẵn là $\mathrm{A}_3^2$.\\
		Số cách chọn $2$ vé cho hai bạn muốn  ngồi ghế bên lẻ là $\mathrm{A}_3^2$. \\
		Còn  lại $2$ vé cho hai bạn còn lại có $2!$ cách. \\
		Vậy số cách chọn là $\mathrm{A}_3^2\cdot  \mathrm{A}_3^2\cdot 2!=72$ cách.
	} 
\end{ex}

% nguồn: MaHoaNH_TracNghiem_Lop10
\begin{ex}%[0D8H2-7]%[ma:3a5bd383]
	Có $4$ cặp vợ chồng được xếp ngồi trên một chiếc ghế dài có $8$ chỗ. Biết rằng mỗi người vợ chỉ ngồi cạnh chồng của mình hoặc ngồi cạnh  một người phụ nữ khác. Hỏi có bao nhiêu cách sắp xếp chỗ ngồi thỏa mãn.
	 \choice
		{\True $816$} 
		{$18$} 
		{$8!$} 
		{$604$} 
	 \loigiai{
		\begin{itemize}
			\item TH1: Chỉ có một cặp vợ chồng ngồi cạnh nhau,  khi đó buộc các bà vợ phải ngồi cùng một bên, các ông chồng ngồi cùng một bên so với cặp vợ chồng đó $\Rightarrow $ có $(2\cdot 3!\cdot 3!)\cdot \mathrm{A}_4^1=288$ (cách xếp).
			\item TH2: Có đúng hai cặp vợ chồng ngồi cạnh nhau suy ra có $2\cdot \mathrm{A}_4^2\cdot 2\cdot 6=288$ (cách  xếp).
			\item TH3: Có đúng ba cặp vợ chồng ngồi cạnh nhau  suy ra có $2\cdot \mathrm{A}_4^3\cdot 2\cdot 2=192$ (cách xếp).
			\item TH4: Tất cả bốn cặp  vợ chồng ngồi cạnh nhau suy ra có $2\cdot \mathrm{A}_4^4=48$ (cách xếp).
		\end{itemize}
		Vậy có tất cả $288+288+192+48=816$ (cách xếp) thỏa mãn yêu cầu  đề bài.
	 }
\end{ex}


%%%%%%%%%%%%%%%% ĐÚNG/SAI %%%%%%%%%%%%%%%%

% nguồn: MaHoaNH_DungSai_Lop10
\begin{ex}%[0D8H2-7]%[ma:fb970ca0]
	Một trường trung học phổ thông có $20$ bạn học sinh tham dự tọa đàm về tháng Thanh niên do  Quận Đoàn tổ chức. Vị trí ngồi của trường là khu vực gồm $4$ hàng ghế, mỗi hàng có $6$ ghế. Xét tính đúng sai của các mệnh đề sau.
	 \choiceTF
		{Có $\mathrm{C}_{20}^6$ cách sắp xếp $6$ bạn ngồi vào hàng ghế  đầu tiên}
		{\True Sau khi sắp xếp  xong hàng ghế đầu tiên, có $\mathrm{A}_{14}^6$ cách sắp xếp $6$ bạn ngồi vào hàng ghế thứ hai}
		{\True Sau khi sắp xếp xong hàng ghế thứ hai, có $\mathrm{A}_8^6$  cách sắp xếp $6$ bạn ngồi vào hàng ghế thứ ba}
		{Sau khi sắp xếp xong hàng ghế thứ ba, có $\mathrm{C}_6^2$ cách sắp  xếp các bạn còn lại ngồi vào hàng ghế cuối cùng}
	\loigiai{ 
		\begin{itemchoice}
			\itemch Mỗi cách chọn $6$ bạn trong $20$ bạn để ngồi vào hàng ghế đầu tiên là một chỉnh hợp chập $6$ của $20$. Vậy có $\mathrm{A}_{20}^6$ cách xếp  $6$ bạn ngồi vào hàng ghế đầu tiên.
			\itemch Mỗi cách chọn $6$ bạn trong $14$ bạn  để ngồi vào hàng ghế thứ hai là một chỉnh hợp chập $6$ của $14$. Vậy có $\mathrm{A}_{14}^6$ cách xếp $6$ bạn ngồi vào hàng ghế thứ hai sau khi sắp xếp xong hàng ghế đầu tiên.
			\itemch Mỗi cách chọn  $6$ bạn trong $8$ bạn để ngồi vào hàng ghế thứ ba là một chỉnh hợp chập $6$ của $8$. Vậy có $\mathrm{A}_8^6$ cách xếp $6$ bạn ngồi vào hàng ghế thứ ba sau khi sắp xếp xong hai hàng ghế đầu.
			\itemch Còn lại $2$ bạn ngồi vào hàng ghế cuối cùng. Mỗi cách chọn $2$ ghế trong $6$ ghế để xếp chỗ ngồi cho $2$ bạn là một chỉnh hợp chập $2$ của 6. Vậy có $\mathrm{A}_6^2$ cách xếp  $2$ bạn còn lại ngồi vào hàng ghế cuối cùng.
		\end{itemchoice}
	 }
\end{ex}

% nguồn: MaHoaNH_DungSai_Lop10
\begin{ex}%[0D8H2-7]%[ma:4f8b03a2]
	An và Bình cùng $7$ bạn khác rủ nhau đi xem bóng đá. Cả $9$ bạn được xếp vào  $9$ ghế theo hàng ngang.
	 \choiceTF
		 {Có $5\,040$ cách xếp chỗ ngồi}
		{\True Có $40\,320$  cách xếp bạn An ngồi chính giữa}
		 {\True Có $80\,640$ cách xếp An và Bình ngồi cạnh nhau}
		{\True Có $282\,240$ cách xếp An và Bình không ngồi  cạnh nhau}
	\loigiai{ 
		\begin{itemchoice}
			\itemch 
			Xếp tùy ý $9$ bạn lên hàng ghé nằm  ngang, ta có $9!=362\,880$ (cách).
			 \itemch
			Xếp bạn An ngồi chính giữa, hoán vị $8$ bạn còn lại ta  có $8!=4\,0320$ (cách).
			 \itemch
			Xếp chỗ cho An và Bình ngồi cạnh nhau (thành nhóm $X$), số cách  xếp trong $X$ là $2!$.\\
			Số cách xếp nhóm $X$ với $7$ người còn lại (ta xem là hoán vị của $8$ phần tử), số cách xếp là  $8!$.\\
			Số cách xếp hàng thỏa mãn là $2!\cdot  8!=80\,640$ (cách).
			\itemch 
			Số cách xếp $9$ bạn vào $9$  chỗ là $9 !$ cách. \\
			Vậy số cách xếp  để An và Bình không ngồi cạnh nhau là $9!-2!\cdot 8!=282\,240$ (cách).
		\end{itemchoice}
	} 
\end{ex}

% nguồn: MaHoaNH_DungSai_Lop10
\begin{ex}%[0D8H2-7]%[ma:b44e4b0d]
	Lớp $10$A có $30$ học sinh  gồm $13$ bạn nữ và $17$ bạn nam, trong đó bạn Khang (nam) làm lớp trưởng.
	 \choiceTF
		{Chọn ra hai bạn gồm một nam  và một nữ tham gia vào Đội cờ đỏ. Số cách chọn là $220$ cách}
		{\True Chọn ra ba bạn trực nhật lớp, trong đó phân công một bạn quét lớp, một bạn quét sân và một bạn lau  bảng. Số cách chọn là $24\,360$}
		 {\True Chọn ra ba bạn tham gia hoạt động thiện nguyện, trong đó phải có lớp trưởng và có ít nhất một nữ. Số cách chọn là $286$}
		{Sắp xếp học sinh để chụp ảnh kỉ yếu trong đó có $14$ bạn đứng hàng trước và $16$ bạn đứng hàng sau.  Số cách sắp xếp là $14!\cdot 16!$}
	 \loigiai{
		\begin{itemchoice}
			 \itemch
			Số cách chọn ra hai bạn gồm một nam và một nữ là $13\cdot 17=221$ cách. 
			\itemch 
			Số cách chọn $3$ bạn từ $30$ bạn và có sắp xếp theo công  việc là $\mathrm{A}_{30}^3=24\,360$.
			 \itemch
			\begin{enumerate}[\bf Trường hợp 1.]
				\item $1$ bạn nữ và $2$ bạn nam (trong đó có bạn Khang lớp trưởng và một bạn nam trong $16$ nam còn lại), số cách chọn $13\cdot  16=208$.
				\item $2$ bạn nữ và bạn Khang lớp trưởng, số  cách chọn $\dfrac{13\cdot 12}{2}\cdot 1=78$.
			\end{enumerate}
			Vậy số cách  chọn là $208+78=286$.
			\itemch 
			Số cách sắp  xếp là $\mathrm{A}_{30}^{14}\cdot 16!$.
		\end{itemchoice}
	} 
\end{ex}

% nguồn: MaHoaNH_DungSai_Lop10
\begin{ex}%[0D8H2-7]%[ma:cf5f5c5a]
	Sắp xếp $5$ bạn học sinh An, Bình, Chi, Dũng, Lan vào một chiếc ghế  dài có $5$ chỗ ngồi.
	\choiceTF 
		{\True Có $120$ cách sắp xếp  tùy ý}
		{Có $16$ cách sắp xếp sao cho bạn Chi luôn ngồi chính  giữa}
		{\True Có $12$ cách sắp xếp sao cho bạn An và  bạn Dũng luôn ngồi ở hai đầu ghế}
		{\True Có $48$ cách sắp xếp sao cho An và Bình ngồi cạnh  nhau}
	 \loigiai{
		\begin{itemchoice}
			\itemch 
			 Số cách sắp xếp $5$ bạn vào ghế dài $5$ chỗ ngồi là $\mathrm{P}_5=5!=120$.
			 \itemch
			Xếp bạn Chi ngồi giữa có  $1$ cách.\\
			Số cách  xếp $4$ bạn sinh An, Bình, Dũng, Lệ vào $4$ chỗ còn lại là một hoán vị của $4$ phần tử nên có có $4!$ cách.\\
			 Vậy có $24$ cách xếp.
			 \itemch
			Xếp An và Dũng ngồi hai đầu  ghế có $2!$ cách xếp.\\
			Số cách xếp $3$ bạn Bình, Chi, Lệ vào $3$ ghế còn lại là một hoán vị của  $3$ phần tử nên có có $3!$ cách.\\
			Vậy có $2!\cdot 3!=12$ cách. 
			\itemch 
			Coi An và Bình là một phần  tử, $3$ học sinh còn lại mỗi bạn là một phần tử.\\
			Xếp $4$ phần  tử này vào ghế dài thì có $4!$ cách.\\
			 Tương ứng với mỗi cách xếp ở trên thì sẽ có $2!$ cách xếp An và Bình ngồi cạnh nhau.\\
			Theo quy tắc nhân ta có  $4!\cdot 2!=48$ cách sắp xếp An và Bình ngồi cạnh nhau.
		\end{itemchoice}
	 }
\end{ex}

% nguồn: MaHoaNH_DungSai_Lop10
\begin{ex}%[0D8H2-7]%[ma:198ed101]
	Có $5$ nam sinh và  $3$ nữ sinh cần được xếp vào một hàng dọc.
	 \choiceTF
		{\True  Có $4\,0320$ cách xếp tùy ý}
		{Có $720$ cách xếp học sinh cùng  giới đứng cạnh nhau}
		{\True Có $4\,320$ cách xếp học sinh  nữ luôn đứng cạnh nhau}
		{\True Có $2\,880$ cách xếp học  sinh nam luôn đứng cạnh nhau}
	\loigiai{ 
		\begin{itemchoice}
			\itemch 
			Số cách xếp  $8$ học sinh theo một hàng dọc $\mathrm{P}_8=8!=4\,0320$ (cách).
			 \itemch
			Gọi $X$ là nhóm $3$ học sinh nữ, $Y$  là nhóm $5$ học sinh nam. \\
			 Số cách xếp trong $X$ là $3!$.\\
			Số  cách xếp trong $Y$ là $5!$. \\
			Số cách hoán  đổi $X$, $Y$ là $2!$.\\
			Vậy số cách xếp thỏa mãn đề bài $3!\cdot 5!\cdot  2!=1\,440$ (cách).
			 \itemch
			Gọi $X$ là nhóm  $3$ học sinh nữ. Khi ấy số cách xếp trong $X$ lad $3!$. \\
			Số cách xếp nhóm $X$ với $5$ học sinh  nam $6 !$.\\
			Vậy  số cách xếp thỏa mãn đề bài $3!\cdot 6!=4\,320$ (cách).
			 \itemch
			Gọi  $X$ là nhóm $5$ học sinh nam. Khi ấy số cách xếp trong $X$ là $5!$. \\
			 Số cách xếp nhóm $X$ với $3$ học sinh nữ $4 !$.\\
			Vậy số cách xếp thỏa mãn  đề bài $5!\cdot 4!=2\,880$ (cách).
		\end{itemchoice}
	} 
\end{ex}

% nguồn: MaHoaNH_DungSai_Lop10
\begin{ex}%[0D8H2-7]%[ma:88614e5c]
	Cho một bàn dài có  $10$ ghế và $10$ học sinh trong đó có $5$ học sinh nữ.
	 \choiceTF
		 {Có $5$ cách xếp $5$ học sinh nữ vào $5$ chỗ ngồi}
		{\True Có $10!$  cách xếp $10$ học sinh vào $10$ ghế}
		{Có $5!\cdot 5!$ cách xếp nam, nữ ngồi  xen kẽ nhau}
		{Có $2\cdot 5!$ cách xếp  học sinh cùng giới ngồi cạnh nhau}
	 \loigiai{
		\begin{itemchoice}
			 \itemch
			Có $5!$ cách xếp $5$ học  sinh nữ vào $5$ chỗ ngồi.
			\itemch 
			Có $10!$ cách  xếp $10$ học sinh vào $10$ ghế.
			 \itemch
			\begin{enumerate}[\bf Trường hợp 1.]
				\item Xếp $5$ học sinh nam ngồi vào vị  trí chẵn có $5!$ cách, sau đó xếp $5$ học sinh nữ vào $5$ vị trí còn lại có $5!$ cách.
				\item Xếp $5$ học sinh nam ngồi ở vị trí lẻ cũng tương  tự như trên.
			\end{enumerate}
			 Vậy có $2\cdot 5!\cdot 5!$ cách xếp nam, nữ ngồi xem kẽ nhau.
			\itemch 
			\begin{itemize}
				\item Xem $5$ nam là $1$ tổ và $5$ nữ là một  tổ, ta có $2$ tổ.
				\item Xếp $2$ tổ ngồi vào bàn  ta có $2!$ cách.
				\item Đổi chỗ $5$ nam cho nhau có $5!$ cách, đổi chỗ $5$ nữ  cho nhau có $5!$ cách.
			\end{itemize}
			Vậy  ta có $2!\cdot 5!\cdot 5!=28\,800$ cách.
		\end{itemchoice}
	 }
\end{ex}

% nguồn: MaHoaNH_DungSai_Lop10
\begin{ex}%[0D8H2-7]%[ma:0ba1626e]
	Xếp $5$ học sinh nam và $2$ học sinh nữ vào  một ghế dài.
	\choiceTF 
		{\True  Có $5\,040$ cách xếp ngẫu nhiên}
		{Có $240$ cách xếp để học sinh cùng giới ngồi cạnh nhau} 
		{Có  $240$ cách xếp để $2$ học sinh nữ ngồi ở $2$ đầu ghế}
		{\True Có $3\,600$ cách xếp để  $2$ học sinh nữ không ngồi cạnh nhau}
	 \loigiai{
		\begin{itemchoice}
			\itemch 
			Số cách xếp ngẫu nhiên $7$ học sinh  không kể nam nữ lên ghế là $\mathrm{P}_7=5\,040$.
			\itemch 
			Các học sinh cùng giới ngồi  cạnh nhau, ta coi các bạn nam là nhóm A, các bạn nữ là nhóm.
			 \itemch
			Xếp $2$ nhóm này lên ghế  có $2!=2$ cách.\\
			 Hoán vị $5$ học sinh nam có $5!=120$ cách.\\
			Hoán  vị $2$ học sinh nữ có $2!=2$ cách.\\
			Vậy số  cách xếp để học sinh cùng giới ngồi cạnh nhau là $2\cdot 120\cdot 2=480$ cách.
			\itemch 
			Xếp $2$ học sinh nữ vào $2$ đầu ghế có  $2!=2$ cách.
			Xếp $5$ học  sinh nam vào $5$ vị trí ở giữa có $5!=120$ cách.
			Vậy số cách xếp để $2$ học sinh nữ ngồi ở $2$ đầu ghế là  $2\cdot 120=240$ cách.\\
			Để $2$ học sinh nữ  ngồi cạnh nhau ta coi $2$ học sinh nữ là nhóm A.\\
			Xếp nhóm $A$ và  $5$ học sinh nam ghế có $6!=720$ cách.\\
			 Hoán vị $2$ học sinh nữ có $2!=2$ cách.\\
			Vậy  số cách xếp để $2$ học sinh nữ ngồi cạnh nhau là $720\cdot 2=1\,440$ cách.\\
			Suy ra xếp $7$ học sinh vào ghế, số cách  xếp để $2$ học sinh nữ không ngồi cạnh nhau là $5\,040-1\,440=3\,600$.
		\end{itemchoice}
	 }
\end{ex}

% nguồn: MaHoaNH_DungSai_Lop10
\begin{ex}%[0D8H2-7]%[ma:03e58618]
	Từ $5$ học sinh không có bạn nào trùng tên nhau trong đó có bạn  Hoa và Hồng. Xếp $5$ bạn đó vào một bàn dài $5$ chỗ.
	 \choiceTF
		{\True  Có $120$ cách xếp tùy ý}
		{Có $24$ cách xếp Hoa  và Hồng ngồi cạnh nhau}
		{\True Có $12$ cách xếp Hoa và Hồng ngồi ở  hai đầu bàn}
		 {\True Có $36$ cách xếp Hoa và Hồng ngồi cách nhau đúng một bạn}
	 \loigiai{
		\begin{itemchoice}
			\itemch 
			Số cách xếp tùy ý $5$ bạn vào bàn dài $5$ chỗ là $\mathrm{P}_5=5!=120$. 
			 \itemch
			Coi $2$ bạn Hoa và Hồng là một phần tử, $3$ học sinh còn  lại mỗi học sinh là một phần tử.\\
			Xếp $4$ phần tử này vào bàn dài thì có $4!$  cách.\\
			Ứng với mỗi cách xếp đó lại có $2!$ cách mà Hoa và Hồng ngồi  cạnh nhau.\\
			Do đó có $4!\cdot  2!=48$ cách xếp Hoa và Hồng ngồi cạnh nhau.
			 \itemch
			Xếp $2$ bạn Hoa và Hồng vào $2$ đầu bàn có $2!$  cách xếp.\\
			Xếp $3$ bạn vào  $3$ vị trí còn lại có $3!$ cách.\\
			Vậy có  $2!\cdot 3!=12$ cách xếp Hoa và Hồng ngồi ở hai đầu bàn.
			 \itemch
			Giả sử bàn dài $5$ chỗ được đánh  số như hình vẽ
			\begin{center}
				\begin{tabular}{|c|c|c|c|c|}
					 \hline
					1 & 2 &  3 & 4 & 5 \\
					 \hline
				\end{tabular}
			\end{center}
			Có $3$ cách chọn vị trí sao cho khi xếp Hoa và Hồng vào thì cách  nhau đúng $1$ bạn đó là các vị trí $1$ và $3$; $2$ và $4$; $3$ và $5$.\\
			Ứng với mỗi vị trí đó xếp $2$ bạn Hoa và Hồng  vào thì có $2!$ cách.\\
			Xếp $3$  bạn vào $3$ vị trí còn lại thì có $3!$ cách.\\
			Do đó có $3\cdot  2!\cdot 3!=36$ cách xếp thỏa mãn yêu cầu bài toán.
		\end{itemchoice}
	} 
\end{ex}

% nguồn: MaHoaNH_DungSai_Lop10
\begin{ex}%[0D8V2-7]%[ma:4342408b]
	Một đoàn tàu nhỏ có $3$ toa khách đỗ  ở sân ga. Có $3$ hành khách không quen biết cùng bước lên tàu. Khi đó, xét tính đúng sai của các mệnh đề sau.
	 \choiceTF
		{Số khả năng khách lên tàu tùy ý  là $9$ khả năng}
		{Số khả  năng $3$ hành khách lên cùng một toa là $1$ khả năng}
		{\True Số khả năng  mỗi khách lên một toa là $6$ khả năng}
		{\True Số khả năng có $2$ hành khách cùng lên một toa, hành khách thứ ba thì  lên toa khác là 18}
	 \loigiai{
		\begin{itemchoice}
			\itemch Khách lên tàu tùy ý nên mỗi khách sẽ có $3$ lựa chọn. Vậy số khả năng thỏa mãn là $3 \cdot 3 \cdot  3=27$.
			\itemch Số khả năng  $3$ hành khách lên cùng một toa là $3$.
			\itemch Số cách chọn  $3$ toa để xếp $3$ hành khách là $\mathrm{A}_3^3=3!=6$.
			\itemch {\bfseries Giai đoạn 1.} Chia $3$ hành khách ra làm hai nhóm X, Y: một nhóm có $2$ người và một nhóm có $1$ người. Số cách thực hiện  là $\mathrm{C}_3^2 \cdot 1$.\\
			{\bfseries Giai đoạn 2.} Chọn $2$ trong  $3$ toa tàu để xếp hai nhóm vào, số cách thực hiện là $\mathrm{A}_3^2$.\\
			Vậy số cách xếp khách  lên tàu thỏa mãn là $\mathrm{C}_3^2 \cdot 1 \cdot \mathrm{A}_3^2=18$.
		\end{itemchoice}
	 }
\end{ex}

% nguồn: MaHoaNH_DungSai_Lop10
\begin{ex}%[0D8V2-7]%[ma:441fab3a]
	Một nhóm $10$ học sinh gồm $5$ học sinh nam trong đó có An và $5$ học  sinh nữ trong đó có Bình được xếp ngồi vào $10$ cái ghế trên một hàng ngang. Xét tính đúng sai của các khẳng định sau.
	 \choiceTF
		{\True Có $120$ cách chọn $3$ học sinh từ nhóm học sinh đó tham gia đại  hội Đoàn trường}
		{\True Có $24$ cách xếp $5$ học sinh nữ vào chiếc ghế dài có $5$ chỗ  ngồi sao cho bạn Bình luôn ngồi chính giữa}
		 {Có $45$ cách chọn ra $2$ học sinh từ nhóm đó để làm tổ trưởng và tổ phó}
		{\True Có $18432$ cách xếp nam và nữ ngồi xen  kẽ, đồng thời An không ngồi cạnh Bình}
	\loigiai{ 
	\begin{itemchoice}
		\itemch Chọn $3$ học sinh từ $10$ học  sinh là $\mathrm{C}_{10}^3=120$.
		\itemch Xếp Bình ngồi giữa có  $1$ cách.\\
		Xếp $4$ bạn còn lại vào $4$ chỗ là  $4!$.\\
		Vậy số cách  là $4!=24$.
		\itemch Chọn $2$ học sinh phân biệt nhận chức vụ khác  nhau $\mathrm{A}_{10}^{2}=90$.
		 \itemch Xếp $5$ nam và $5$ nữ xen kẽ có $2\cdot 5!\cdot 5!$ cách.\\
		 Xếp An và Bình ngồi cạnh nhau.
		\begin{itemize}
			\item Chọn chỗ cho An và Bình ngồi cùng nhau  có $9$ cách.
			 \item An và Bình hoán đổi vị trí $2!$.
			\item Xếp $4$ nam và $4$ nữ  còn lại xen kẽ là $4!\cdot4!$ cách.
		\end{itemize}
		Vậy số cách thỏa An, Bình  không ngồi cạnh nhau là $2\cdot5!\cdot 5!-9\cdot2!\cdot4!\cdot4!=18432$.
	\end{itemchoice}
	} 
\end{ex}

% nguồn: MaHoaNH_DungSai_Lop10
\begin{ex}%[0D8V2-7]%[ma:307f36cc]
	Xếp $6$ bạn trong  đó có Bình và An thành một hàng dọc.
	 \choiceTF
		{\True Có $6!$ cách xếp bất kì} 
		{\True Có $4!$  cách xếp sao cho Bình hoặc An đứng đầu hàng}
		{\True Có $2\cdot 5!$ cách xếp sao cho  Bình và An đứng cạnh nhau}
		{\True Có $4!\cdot\mathrm{A}_5^2$ cách xếp Bình và  An không đứng cạnh nhau}
	\loigiai{ 
		\begin{itemchoice}
			 \itemch
			Số các  xếp $6$ bạn thành một hàng dọc là $6!$.
			 \itemch
			Bình hoặc An đứng đầu hàng có  $2$ cách.\\
			 Xếp $4$ bạn còn lại có $4!$ cách.\\
			Vậy  số cách xếp $6$ bạn sao cho Bình hoặc An đứng đầu hàng là $2\cdot 4!$.
			 \itemch
			Xếp Bình và An đứng cạnh  nhau có $2$ cách.\\
			Xếp  Bình, An và $4$ bạn còn lại có $5!$ cách.\\
			Vậy có $2\cdot  5!$ cách xếp sao cho Bình và An đứng cạnh nhau.
			\itemch 
			Xếp $4$  bạn còn lại có $4!$ cách.\\
			Xếp Bình và An vào $2$ trong $5$  khe trống có $\mathrm{A}_5^2$ cách.\\
			Vậy có $4!\cdot\mathrm{A}_5^2$ cách xếp  Bình và An không đứng cạnh nhau.
		\end{itemchoice}
	 }
\end{ex}

% nguồn: MaHoaNH_DungSai_Lop10
\begin{ex}%[0D8V2-7]%[ma:01182374]
	 Xếp $7$ bạn nam và $6$ bạn nữ thành một hàng ngang.
	\choiceTF 
		{\True Có $13!$ cách sắp xếp các  bạn thành $1$ hàng ngang}
		{Nam nữ đứng xen kẻ thì có $7!\cdot 6!$ cách sắp xếp} 
		{Nữ luôn đứng cạnh nhau thì có $3\,628\,800$ cách sắp  xếp}
		{\True Có $7!\cdot \mathrm{A}_8^6$ cách  xếp không có $2$ nữ nào đứng cạnh nhau}
	 \loigiai{
		\begin{itemchoice}
			\itemch 
			Có $13$ bạn nên có $13!$ cách xếp $13$  bạn thành $1$ hàng ngang.
			 \itemch
			\begin{enumerate}[\bf Trường hợp 1.]
				\item Bạn nam đứng đầu có $7$ cách chọn, kế đến là bạn nữ có $6$ cách chọn, kế đến là bạn nam có $6$ cách chọn, kế đến là $1$ bạn nữ có $5$ cách chọn; cuối cùng  xếp $1$ bạn nữ có $1$ cách chọn. \\
				Suy ra tổng số cách xếp  $7!\cdot 6!$ cách.
				\item Bạn nữ đứng đầu, xếp hoàn toàn tương tự như trường  hợp 1.\\
				Suy ra tổng số cách sếp  của trường hợp này là $6!\cdot 7!$.
			\end{enumerate}
			Kết luận theo quy tắc cộng tổng số cách xếp nam nữ xen kẽ nhau là \begin{align*}
			7!\cdot 6!+6!\cdot 7!=2\cdot 7!\cdot 6!. 
		\end{align*}
		\itemch 
		Gọi  nhóm bạn nữ là nhóm X. Số cách xếp $7$ bạn nam và X là $8!$ cách.\\
		Ứng với mỗi cách xếp  trên có $6!$ cách xếp $6$ bạn nữ trong nhóm X.\\
		 Theo quy tắc nhân có $8!\cdot 6!=29\,030\,400$ cách xếp.
		\itemch 
		Bước đầu tiên xếp $7$ bạn nam đứng kề nhau có  $7!$ cách xếp. \\
		Để các bạn nữ không đứng kế nhau ta xen các bạn nữ vào giữa các bạn nam, giữa $7$ bạn  nam có $6$ vị trí và thêm $2$ vị trí đầu và cuối, tổng cộng có $8$ vị trí để xếp $6$ bạn nữ.\\
		Chọn $6$ vị trí trong $8$ vị trí để xếp các  bạn nam, có $\mathrm{A}_8^6$ cách.\\
		Theo quy tắc nhân có $7!\cdot  \mathrm{A}_8^6$ cách xếp thỏa yêu cầu bài toán.
	\end{itemchoice}
 }
\end{ex}

% nguồn: MaHoaNH_DungSai_Lop10
\begin{ex}%[0D8T2-7]%[ma:0a3879f7]
	Có hai dãy ghế, mỗi dãy  $5$ ghế. Xếp $5$ nam, $5$ nữ vào hai dãy ghế trên. Chọn khẳng định đúng trong các khẳng định sau.
	 \choiceTF
		{Có $14400$ cách xếp nam,  nữ ngồi tùy ý}
		{\True Có $3628800$  cách xếp nam, nữ ngồi tùy ý}
		{Có $14400$ cách xếp  nam một dãy ghế, một dãy ghế}
		{\True Có $28800$ cách xếp nam một dãy ghế, nữ  một dãy ghế}
	 \loigiai{
		\begin{itemize}
			\item Xếp nam,  nữ ngồi tùy ý có $10!=3628800$.
			\item Chọn $1$ dãy xếp nam ngồi có $2$ cách, xếp $5$ nam vào dãy đã chọn có $5!$ cách, xếp nữ vào dãy ghế  còn lại có $5!$ cách. Vậy có $2\cdot 5! \cdot 5!=28800$ cách xếp nam một dãy ghế, nữ một dãy ghế.
		\end{itemize}
	 }
\end{ex}

% nguồn: MaHoaNH_DungSai_Lop10
\begin{ex}%[0D8T2-7]%[ma:98031a99]
	Trên một kệ sách dài có $ 5 $ quyển sách  Toán, $ 4 $ quyển sách Lí, $3$ quyển sách Văn. Các quyển sách đều khác nhau. Chọn khẳng định đúng trong các khẳng định sau.
	 \choiceTF
		{Có $3!\cdot 5!\cdot  4!$ cách sắp xếp tùy ý các quyển sách trên}
		{\True Có $12!$ cách sắp  xếp tùy ý các quyển sách trên}
		{Có $3!\cdot 5!\cdot 4!$ cách sắp xếp các quyển sách trên  theo từng môn}
		{\True Có $103680$ cách sắp xếp các quyển  sách trên theo từng môn}
	 \loigiai{
		\begin{itemize}
			\item Xếp $12$ quyển sách trên có $12!$ cách  xếp tùy ý.
			\item Xem  nhóm $5$ sách Toán thành khối $A$, $4$ sách Lí thành khối $B$, $3$ sách Văn thành khối $C$. Xem đây là $3$ hoán vị của $3$ phần tử $A$, $B$, $C$. Suy ra, có $3!$ cách xếp $A$, $B$, $C$. Xếp $5$ sách Toán trong khối $A$ có $5!$ cách, xếp $4$ sách Lí trong khối $B$ có $4!$ cách, xếp $3$ sách Văn trong khối $C$ có $3!$ cách. Theo quy tắc nhân, có $3!\cdot 5!\cdot 4!\cdot 3!=103680$ cách xếp.
		\end{itemize}
	} 
\end{ex}


%%%%%%%%%%%%%%%% TRẢ LỜI NGẮN %%%%%%%%%%%%%%%%

% nguồn: MaHoaNH_TraLoiNgan_Lop10
\begin{bt}%[0D8H2-7]%[ma:5a004a77]
	Có $9$ cuốn sách khác nhau gồm $3$ cuốn sách Lý, $4$ cuốn sách Sinh, $2$ cuốn sách Địa. Hỏi có bao nhiêu cách sắp xếp  các cuốn sách trên vào giá sách hàng ngang nếu các cuốn sách cùng môn học đứng cạnh nhau?
	 \shortans{1728}
	 \loigiai{
		 Gọi L là nhóm $3$ sách lý, S là nhóm $4$ sách sinh, Đ là nhóm $2$ sách địa.\\
		Số cách xếp trong L là $3!$; số cách xếp trong S là $4!$; số cách xếp trong Đ là $2!$; số cách xếp L, S, Đ với nhau  là $3!$.\\
	Vậy số cách xếp thỏa  mãn đề bài là $3!\cdot 4!\cdot 2!\cdot 3!=1\ 728$ (cách).}
\end{bt}

% nguồn: MaHoaNH_TraLoiNgan_Lop10
\begin{bt}%[0D8H2-7]%[ma:e9888191]
	Có bao nhiêu cách sắp xếp $2$ nữ sinh, $2$ nam sinh thành một hàng ngang sao cho các bạn nam và các bạn nữ ngồi xen kẽ? 
	\shortans{8} 
	\loigiai{ 
		 Đánh số hàng ngang là $1$; $2$; $3$; $4$.\\
		Xếp $2$ bạn nam ngồi ở vị trí số $1$, $3$ và hoán  vị $2$ bạn này ta có $2!$ (cách).\\
		Xếp $2$ bạn nữ  ngồi ở vị trí số $2$, $4$ và hoán vị $2$ bạn này ta có $2!$ (cách).\\
		Xem $2$ bạn nam ngồi ở vị số $1$, $3$ là một phần tử với $2$ bạn nữ ngồi ở vị  trí số $2$, $4$ là một phần tử hoán vị vị trí ta có $2!$ cách.\\
		Vậy có $8$ cách sắp  xếp thỏa mãn yêu cầu bài toán.
	 }
\end{bt}


%%%%%%%%%%%%%%%% TỰ LUẬN %%%%%%%%%%%%%%%%

% nguồn: MaHoaNH_TuLuan_Lop10
\begin{bt}%[0D8H2-7]%[ma:f54cabbf]
	Có bao nhiêu cách sắp chỗ ngồi cho $6$ người vào  $6$ ghế xếp thành một dãy?
	 \loigiai{
		Mỗi một cách sắp xếp là  một hoán vị của $6$ phần tử.\\
		Do đó có $\mathrm{P}_6 = 6! = 720$ cách xếp. 
	} 
\end{bt}

% nguồn: MaHoaNH_TuLuan_Lop10
\begin{bt}%[0D8H2-7]%[ma:73baaeac]
	Có bao nhiêu cách xếp khác nhau cho $4$ người ngồi vào  $6$ chỗ trên một ghế dài?
	\loigiai{ 
		Mỗi cách xếp khác nhau cho $4$ người ngồi vào $6$ chỗ trên một ghế dài là một chỉnh hợp chập $4$ của $6$ phần  tử.\\
		 Vậy có $\mathrm{A}_6^4 = 360$ cách.
	} 
\end{bt}

% nguồn: MaHoaNH_TuLuan_Lop10
\begin{bt}%[0D8H2-7]%[ma:84f128ef]
	Có bao nhiêu cách  sắp xếp cho $5$ học sinh nam và $3$ học sinh nữ ngồi quanh một bàn tròn sao cho không có hai học sinh nữ nào cạnh nhau? (Nếu có hai cách sắp xếp mà cách xếp này quay quanh vòng tròn được cách sắp xếp kia thì ta coi chỉ là một cách sắp xếp).
	 \loigiai{
		Giả sử đã xếp chỗ cho $5$ học sinh nam. Vì ba học sinh nữ không ngồi cạnh nhau nên họ được chọn $3$ trong $5$ vị trí xen kẽ giữa các học sinh nam, số cách chọn là $\mathrm{A}_5^3$.  Vì hai cách xếp vị trí cho $8$ người với cùng một thứ tự quay quanh bàn tròn được coi là một nên ta có thể chọn trước vị trí cho một học sinh nam nào đó, số hoán vị của $4$ học sinh nam còn lại vào các vị trí là $4!$.\\
		Theo quy tắc nhân, số khả năng phải tìm  là: $4! \cdot \mathrm{A}_5^3 = 1440$ (cách).\\
		Vậy số  cách xếp thỏa mãn là: $1\,440$ cách.
	 }
\end{bt}

% nguồn: MaHoaNH_TuLuan_Lop10
\begin{bt}%[0D8H2-7]%[ma:6a72154b]
	Có bao nhiêu cách sắp xếp $5$ viên bi  đỏ khác nhau và $5$ viên bi đen khác nhau thành một dãy sao cho hai viên bi cùng màu không xếp cạnh nhau?
	\loigiai{ 
		Sắp xếp $5$ bi đỏ,  có $5!$ cách.\\
		Chọn vị trí để sắp xếp bi đen xen giữa các bi đỏ, có $2$ cách (bi đen đứng đầu hoặc bi  đỏ đứng đầu).\\
		Sắp xếp $5$ bi đen vào  vị trí đã chọn, có $5!$ cách.\\
		Vậy số cách sắp xếp thỏa  mãn yêu cầu bài toán là $5! \cdot 2 \cdot 5! = 28800$ cách.
	 }
\end{bt}

% nguồn: MaHoaNH_TuLuan_Lop10
\begin{bt}%[0D8H2-7]%[ma:b26d5df0]
	Xếp $6$ bạn học sinh nam và $3$ bạn học sinh nữ $A$, $B$, $C$ ngồi trên một hàng ngang có $9$ ghế sao cho mỗi ghế có đúng một học sinh. Số cách xếp chỗ ngồi cho $9$ học sinh đó sao cho mỗi bạn nữ  ngồi giữa hai học sinh nam là.
	\loigiai{ 
		Xếp $6$ học sinh nam  có $6!$ cách xếp.\\
		Giữa $6$ học sinh nam có $5$  khoảng trống.\\
		Chọn $3$ khoảng trống trong $5$ khoảng trống trên và xếp $3$ học sinh  nữ $A$, $B$, $C$ vào có: $\mathrm{A}_5^3$ cách.\\
		Vậy số cách xếp thỏa mãn yêu cầu là: $6! \cdot \mathrm{A}_5^3 =  43\,200$ cách.
	} 
\end{bt}

% nguồn: MaHoaNH_TuLuan_Lop10
\begin{bt}%[0D8H2-7]%[ma:1dbe069c]
	Có thể lập ra bao nhiêu số điện thoại di động có $10$ chữ số bắt đầu là  $0908$, các chữ số còn lại khác nhau, khác với $4$ chữ số đầu và phải có mặt chữ số $6$.
	\loigiai{ 
		Gọi số điện thoại có dạng $\overline{0908abcdef}$.\\ 
		Chọn $1$ vị trí trong $6$ vị trí $\overline{abcdef}$ để xếp chữ số $6$ có $6$  cách chọn.\\
		Chọn  $5$ chữ số trong $6$ chữ số là $\{1; 2; 3; 4;, 5; 7\}$ để xếp vào $5$ vị trí còn lại, có $\mathrm{A}_6^5$ cách.\\
		Kết luận có $6  \cdot \mathrm{A}_6^5 = 4320$ số điện thoại thỏa yêu cầu.
	 }
\end{bt}

% nguồn: MaHoaNH_TuLuan_Lop10
\begin{bt}%[0D8H2-7]%[ma:2df2f8ca]
	Có hai dãy ghế, mỗi dãy $5$ ghế. Xếp $5$ nam, $5$ nữ vào $2$ dãy ghế trên có bao  nhiêu cách nếu:
	\begin{enumerate}
		\item Nam và nữ được xếp tùy  ý.
		\item Nam  $1$ dãy ghế, nữ $1$ dãy ghế.
	\end{enumerate}
	\loigiai{ 
		\begin{enumerate}
			\item Mỗi cách xếp $5$ nam  và $5$ nữ vào hai dãy ghế một cách tùy ý là một hoán vị của $10$ người.\\
			Vậy có $10! =  3628800$ cách xếp.
			\item Chọn $1$ dãy  để xếp nam ngồi vào có $2$ cách; xếp $5$ nam vào dãy ghế đã chọn có $5!$ cách; xếp $5$ nữ vào dãy ghế còn lại có $5!$ cách.\\
			Vậy có tất cả là $2 \cdot 5! \cdot 5!$ cách xếp thỏa điều kiện bài  toán.
		\end{enumerate}
	 }
\end{bt}

% nguồn: MaHoaNH_TuLuan_Lop10
\begin{bt}%[0D8H2-7]%[ma:bbd03e19]
	Có bao nhiêu cách xếp $5$ bạn học sinh A, B, C,  D, E vào $5$ ghế dài sao cho:
	\begin{enumerate}[a)]
		\item Bạn C ngồi  chính giữa?
		\item Hai bạn A và E ngồi  ở hai đầu ghế?
	\end{enumerate}
	 \loigiai{
		\begin{enumerate}[a)]
			\item Xếp bạn C ngồi ở chính giữa: có  $1$ cách xếp.\\
			Xếp $4$ bạn còn lại vào $4$ vị trí còn lại: có $\mathrm{P}_4 = 4!  = 24$ cách xếp.\\
			Theo quy tắc nhân, có $1 \cdot 24 = 24$ cách xếp. 
			\item Xếp $2$ bạn A và E ngồi  ở hai đầu ghế: có $2! = 2$ cách xếp.\\
			Xếp $3$ bạn còn lại vào $3$ vị trí  còn lại: có $3! = 6$ cách xếp.\\
			Theo quy tắc nhân, có $2! \cdot 3! = 2 \cdot 6 = 24$ cách  xếp.
		\end{enumerate}
	 }
\end{bt}

% nguồn: MaHoaNH_TuLuan_Lop10
\begin{bt}%[0D8H2-7]%[ma:be017e4b]
	Cho một bàn dài có $10$ ghế và $10$ học sinh trong đó có $5$ học sinh nữ. Hỏi có bao  nhiêu cách sắp chỗ ngồi cho $10$ học sinh sao cho:
	\begin{enumerate}
		\item Nam, nữ ngồi xen kẽ  nhau?
		\item Những học sinh cùng  giới thì ngồi cạnh nhau?
	\end{enumerate}
	\loigiai{ 
		\begin{enumerate}
			\item Cách 1: Xếp $5$ học sinh nam ngồi vào vị trí chẵn  có $5!$ cách, sau đó xếp $5$ học sinh nữ vào $5$ vị trí còn lại có $5!$ cách $\Rightarrow$ có $5! \cdot 5!$ cách.\\
			Cách 2: Xếp $5$ học sinh nam ngồi vào vị trí lẻ có $5!$ cách, sau đó xếp $5$ học sinh nữ vào $5$ vị trí còn lại  có $5!$ cách $\Rightarrow$ có $5! \cdot 5!$ cách.\\
			Vậy tất cả có  $2 \cdot 5! \cdot 5! = 28800$ cách.
			\item Xem $5$ nam là $1$ tổ và $5$ nữ là một tổ, ta có $2$ tổ. Xếp  $2$ tổ ngồi vào bàn ta có $2!$ cách.\\
			Đổi chỗ $5$  nam cho nhau có $5!$ cách, đổi chỗ $5$ nữ cho nhau có $5!$ cách.\\
			Vậy ta có $2! \cdot 5! \cdot  5! = 28800$ cách.
		\end{enumerate}
	} 
\end{bt}

% nguồn: MaHoaNH_TuLuan_Lop10
\begin{bt}%[0D8H2-7]%[ma:a41ef0b7]
	Một trường trung học phổ thông có $4$ học sinh giỏi khối $12$, có $5$ học sinh giỏi khối $11$, có $6$ học sinh giỏi khối $10$. Hỏi  có bao nhiêu cách sắp xếp $15$ học sinh trên thành một hàng ngang để đón đoàn đại biểu, nếu:
	\begin{enumerate}
		\item Các học sinh được  xếp bất kỳ.
		\item Các học sinh trong  cùng một khối phải đứng kề nhau.
	\end{enumerate}
	 \loigiai{
		\begin{enumerate}
			\item Mỗi cách sắp xếp $15$ học  sinh thành một hàng ngang là một hoán vị của $15$ phần tử.\\
			Vậy có $15!$ cách xếp $15$ học sinh thành một  hàng ngang.
			\item Xếp các khối có  $3!$ cách xếp.\\
			Xếp các bạn trong khối  $12$ có $4!$ cách.\\
			 Xếp các bạn trong khối $11$ có $5!$ cách.\\
			Xếp các bạn trong khối  $10$ có $6!$ cách.\\
			Theo quy tắc nhân có $3! \cdot 4! \cdot 5! \cdot 6! =  12441600$ cách xếp thỏa yêu cầu.
		\end{enumerate}
	 }
\end{bt}

% nguồn: MaHoaNH_TuLuan_Lop10
\begin{bt}%[0D8H2-7]%[ma:7f7930ba]
	Một chồng sách gồm $4$ quyển sách Toán khác nhau, $3$ quyển sách Vật Lý khác nhau, $5$  quyển sách Hóa Học khác nhau. Hỏi có bao nhiêu cách xếp các quyển sách trên thành một hàng ngang sao cho
	\begin{enumerate}[a)]
		 \item Các quyển sách cùng môn thì đứng cạnh nhau.
		\item  Các quyển sách Toán đứng gần nhau.
	\end{enumerate}
	 \loigiai{
		\begin{enumerate}[a)]
			\item Xếp $4$ quyển sách Toán thành một nhóm đứng gần nhau có $\mathrm{P}_4 = 4!  = 24$ cách xếp.\\
			Xếp $3$ quyển sách Vật Lý thành một nhóm gần  nhau có $\mathrm{P}_3 = 3! = 6$ cách xếp.\\
			Xếp $5$ quyển sách Hóa Học thành một nhóm gần nhau có $\mathrm{P}_5 = 5! = 120$ cách xếp.\\ 
			Xếp $3$ nhóm sách trên lên giá sách có  $\mathrm{P}_3 = 3! = 6$ cách xếp.\\
			Vậy có $24 \cdot 6  \cdot 120 \cdot 6 = 103680$ cách xếp các cuốn sách cùng môn thì đứng cạnh nhau.
			\item Xếp $4$ quyển sách Toán thành một nhóm đứng gần nhau có $\mathrm{P}_4 = 4!  = 24$ cách xếp.\\
			Coi nhóm  sách Toán là một quyển sách lớn, xếp quyển sách lớn đó và $8$ quyển sách còn lại có $\mathrm{P}_9 = 9!$ cách xếp.\\
			Vậy có $24  \cdot 9! = 8709120$ cách xếp các cuốn sách Toán đứng gần nhau.
		\end{enumerate}
	 }
\end{bt}

% nguồn: MaHoaNH_TuLuan_Lop10
\begin{bt}%[0D8H2-7]%[ma:f8399d0d]
	Trên giá sách dài có $5$ quyển sách Toán, $4$ quyển  sách Văn và $3$ quyển sách tiếng Anh. Các quyển sách đều khác nhau. Hỏi có bao nhiêu cách sắp xếp các quyển sách trên sao cho:
	\begin{enumerate}[a)]
		\item Các quyển sách  xếp một cách tùy ý?
		\item Các quyển sách xếp  theo từng môn liền nhau?
		\item Các quyển sách xếp theo từng môn và sách Toán xếp  ở giữa?
	\end{enumerate}
	\loigiai{ 
		\begin{enumerate}[a)]
			\item Trên giá sách có tổng cộng: $5 +  4 + 3 = 12$ quyển sách.\\
			 Mỗi một cách xếp tùy ý các quyển sách trên giá là một hoán vị của $12$ phần tử.\\
			Vậy có $\mathrm{P}_{12} = 12!  = 479001600$ cách xếp.
			\item Xem mỗi loại sách là một khối thống nhất (\lq\lq buộc\rq\rq\, mỗi loại thành $1$ bó), ta có $3!$ cách xếp $3$  khối này. Có $5!$ cách xếp sách Toán, có $4!$ cách xếp sách Văn, có $3!$ cách xếp sách Tiếng Anh.\\
			Theo quy tắc nhân, có  $3! \cdot 5! \cdot 4! \cdot 3! = 103680$ cách xếp.
			\item Xem mỗi loại sách là  một khối thống nhất (như trên), ta có $2!$ cách xếp $2$ môn còn lại ở hai bên sách Toán.\\
			Ứng với mỗi cách, có $5!$ cách  xếp sách Toán, có $4!$ cách xếp sách Văn, có $3!$ cách xếp sách Tiếng Anh.\\
			Theo quy tắc nhân, có $2! \cdot  5! \cdot 4! \cdot 3! = 34560$ cách xếp.
		\end{enumerate}
	 }
\end{bt}

% nguồn: MaHoaNH_TuLuan_Lop10
\begin{bt}%[0D8H2-7]%[ma:eba00de5]
	 Xếp $5$ bạn nam và $5$ bạn nữ thành một hàng dọc. Hỏi có bao nhiêu cách xếp
	\begin{enumerate}
		\item Nam  nữ đứng xen kẽ.
		 \item Nữ luôn đứng cạnh nhau.
		\item Không có $2$ nam  nào đứng cạnh nhau.
	\end{enumerate}
	\loigiai{ 
		\begin{enumerate}
			\item Trường hợp 1: Bạn nam đứng đầu có $5$ cách chọn, kế đến là bạn nữ có $5$ cách chọn, kế đến là bạn nam có $4$ cách chọn, kế đến là $1$ bạn nữ  có $4$ cách chọn, ... cuối cùng xếp $1$ bạn nữ có $1$ cách chọn.\\
			Suy ra tổng số  cách xếp $5! \cdot 5!$ cách.\\
			Trường hợp 2: Bạn nữ đứng đầu, xếp hoàn toàn tương tự như  trường hợp 1, suy ra tổng số cách xếp của trường hợp này là $5! \cdot 5!$.\\
			Kết luận theo quy tắc cộng tổng số cách xếp nam  nữ xen kẽ nhau là $5! \cdot 5! + 5! \cdot 5! = 28800$ cách.
			\item Gọi nhóm bạn nữ là  nhóm $X$. Số cách xếp $5$ bạn nam và $X$ là $6!$ cách.\\
			Ứng với mỗi cách xếp trên có $5!$ cách xếp $5$ bạn nữ  trong nhóm $X$.\\
			Theo quy tắc nhân có $6!  \cdot 5! = 86400$ cách xếp.
			\item Bước đầu tiên xếp $5$ bạn nữ đứng cạnh nhau có $5!$ cách xếp. Để các bạn nam không đứng cạnh nhau, ta xen các bạn nam vào giữa  các bạn nữ.\\
			Giữa $5$ bạn nữ có $4$ vị trí và thêm $2$ vị trí đầu và cuối, tổng cộng có $6$  vị trí để xếp $5$ bạn nam. Chọn $5$ vị trí trong $6$ vị trí để xếp các bạn nam có $\mathrm{A}_6^5$ cách.\\
			Theo quy tắc nhân  có $5! \cdot \mathrm{A}_6^5 = 86400$ cách xếp thỏa yêu cầu bài toán.
		\end{enumerate}
	 }
\end{bt}

% nguồn: MaHoaNH_TuLuan_Lop10
\begin{bt}%[0D8V2-7]%[ma:78b82102]
	Vợ chồng ông An cùng $4$ người con đang lên máy  bay theo một hàng dọc. Có bao nhiêu cách xếp hàng khác nhau sao cho ông An hay vợ ông An đứng ở đầu hoặc cuối hàng?
	 \loigiai{
		Sắp $6$ người vào $6$ vị trí theo hàng dọc có $6!$ cách  sắp xếp.\\
		Sắp vợ chồng ông An vào $2$ trong $4$ vị trí (trừ vị trí đầu và cuối hàng) có $\mathrm{A}_4^2$ cách.\\ 
		Sắp $4$ người con  vào $4$ vị trí còn lại có $4!$ cách.\\
	Vậy  có $6!-\mathrm{A}_4^2 \cdot 4!=432$ cách sắp xếp thỏa mãn yêu cầu đề bài.}
\end{bt}

% nguồn: MaHoaNH_TuLuan_Lop10
\begin{bt}%[0D8V2-7]%[ma:7103543e]
	An và Bình cùng $7$ bạn  khác rủ nhau đi xem bóng đá. Cả $9$ bạn được xếp vào $9$ ghế theo hàng ngang. Khi đó, xét tính đúng sai của các mệnh đề sau.
	\begin{enumerate}
		\item Có bao nhiêu cách xếp để An và Bình ngồi $2$ đầu  dãy ghế?
		\item Có bao nhiêu cách xếp An và Bình không ngồi  cạnh nhau?
	\end{enumerate}
	\loigiai{ 
		\begin{enumerate}
			\item Số cách xếp  để An, Bình ngồi $2$ đầu dãy ghế là $2!\cdot 7!=10\ 080$.
			\item Số cách xếp $9$ bạn vào $9$  chỗ là $9!$ cách.\\
			Xếp chỗ cho An và Bình ngồi cạnh nhau (thành nhóm  $X$), số cách xếp trong $X$ là $2!$.\\
			Số cách xếp nhóm $X$ với $7$ người còn lại (ta xem là hoán vị của $8$  phần tử) là $8 !$.\\
			Suy ra số cách xếp hàng  sao cho An và Bình ngồi cạnh nhau là $2!\cdot 8!=80\ 640$ (cách).\\
			Vậy số cách xếp để An và Bình không ngồi cạnh nhau  là $9!-2!\cdot 8!=28\ 2240$ (cách).
		\end{enumerate}
	} 
\end{bt}
